% ============================================================
% Prólogo del traductor — texto nuevo, no está en el original
% ============================================================
\chapter*{Prólogo del traductor}
\addcontentsline{toc}{chapter}{Prólogo del traductor}
\markboth{Prólogo del traductor}{Prólogo del traductor}

Esta edición traduce al español latino, por primera vez de manera completa
y con el mismo rigor matemático, los dos volúmenes de las lecciones que
Henri Poincaré profesó en la Sorbona durante el curso 1887-1888 y que
Georges Carré publicó en París bajo el título \emph{Théorie mathématique
de la lumière}. El presente volumen corresponde al primero de los dos
tomos originales (1889); el segundo (1892), dedicado a nuevos estudios
sobre la difracción y a la teoría de la dispersión de Helmholtz, se
abordará una vez concluido este.

El propósito de esta edición no es solo traducir, sino presentar el texto
como un libro guía de física matemática apto para un estudiante de
pregrado. Para ello se han tomado tres decisiones editoriales, visibles a
lo largo de todo el volumen:

\begin{enumerate}
  \item Donde Poincaré o sus redactores dan un paso algebraico por evidente
  (\emph{``on obtient immédiatement''}, \emph{``il vient''}, \emph{``on
  démontre''}), el cálculo se desarrolla explícitamente, paso a paso, sin
  omitir nada. Esta ampliación se escribe en el mismo registro formal e
  impersonal del resto del texto, tejida en el lugar exacto donde el
  original es breve, para que se lea como una continuación natural de su
  propia exposición y no como un aparato pedagógico añadido aparte.
  \item Todas las figuras del libro se rehicieron como ilustraciones
  vectoriales, con una paleta de color propia —azul y rojo— que se usa en
  todo el volumen para las portadillas de capítulo, las figuras y las
  cajas que destacan los enunciados verdaderamente centrales: hipótesis
  fundamentales, teoremas y resultados culminantes con nombre propio. Esas
  cajas nunca contienen una demostración, solo el enunciado.
  \item Se conserva en todo momento la numeración original de las
  ecuaciones de Poincaré, incluidos sus frecuentes reinicios dentro de un
  mismo capítulo. Donde el cotejo con el original reveló una errata
  genuina de la edición de 1889 —ya fuera una ya señalada en la fe de
  erratas del propio libro, ya una no registrada en ella—, se documenta en
  nota al pie sin alterar en silencio el razonamiento de Poincaré.
\end{enumerate}

Toda la traducción se verificó ecuación por ecuación contra el texto
francés original. El lector encontrará, dispersas en notas al pie, las
contadas ocasiones en que esa verificación encontró una inconsistencia
real del propio original —casi siempre un error de imprenta— y la forma en
que se resolvió.

\bigskip
\noindent\textit{El traductor}
